\section*{Resumen}
\addcontentsline{toc}{section}{Resumen}

Este trabajo presenta un mapeo sistem\'atico del flujo de trabajo de \speckit, un toolkit de c\'odigo abierto para \sdd{} desarrollado por GitHub, contra un flujo profesional de ingenier\'ia de requisitos basado en normas. El marco normativo de referencia se compone del \ireb{} CPRE Foundation Level Handbook, la norma ISO/IEC/IEEE 29148 y la \textit{INCOSE Guide for Writing Requirements}.

El estudio documenta el flujo real de \speckit{} (sus comandos, artefactos y mecanismos de calidad) a partir de sus fuentes oficiales, y lo contrasta con las actividades fundamentales de la ingenier\'ia de requisitos ---elicitaci\'on, documentaci\'on, validaci\'on y gesti\'on--- as\'i como con los criterios de calidad definidos por las normas de referencia. Para fundamentar el mapeo con evidencia emp\'irica, se construye un corpus de requisitos reales extra\'idos de \emph{changelogs} de seis proyectos \emph{open source} representativos (Appwrite, Authentik, Cal.com, Directus, Medusa y n8n), se formalizan siguiendo una plantilla con atributos IREB, y se introducen en \speckit{} para observar el comportamiento del pipeline \sdd.

El resultado principal es un mapa de alineaci\'on que identifica, para cada actividad y criterio normativo, si \speckit{} lo cubre, lo cubre parcialmente o no lo cubre, respaldado por evidencia de los casos de estudio. El trabajo concluye con recomendaciones para aproximar el flujo \speckit{} a un proceso de ingenier\'ia de requisitos basado en normas, y se identifican las limitaciones inherentes al enfoque de generaci\'on de c\'odigo dirigido por especificaciones.

\vspace{0.5cm}
\noindent\textbf{Palabras clave:} Ingenier\'ia de requisitos, IREB, SpecKit, Spec-Driven Development, ISO 29148, mapeo de procesos, generaci\'on de c\'odigo con IA.
